\documentclass[11pt]{article}
\usepackage[margin=0.86in]{geometry}
\usepackage{amsmath,amssymb,booktabs,tabularx,array,graphicx,xcolor,hyperref}

\hypersetup{
  colorlinks=true,
  linkcolor=blue!55!black,
  citecolor=blue!55!black,
  urlcolor=blue!55!black
}
\newcolumntype{Y}{>{\raggedright\arraybackslash}X}
\newcommand{\microir}{\textsc{micro-ir}}
\newcommand{\cstate}{\mathcal{S}}
\newcommand{\stagebox}[2]{\fbox{\parbox[c][#1][c]{0.93\linewidth}{#2}}}
\setlength{\emergencystretch}{1.5em}

\title{Heterogeneous Cognitive Transformers:\\
Computational and Epistemic Coprocessors}
\author{Working Draft}
\date{August 2026}

\begin{document}
\maketitle

\begin{abstract}
Calculators, formal solvers, databases, and search engines are often grouped
under tool use, although they repair different deficits and need not share an
optimal interface.  We present an evidence-aware position and research program
in which a Transformer remains the learned semantic and generative substrate
while computational coprocessors address competence deficits and epistemic
coprocessors address knowledge deficits.  A common runtime can recognize an
emerging need, normalize it into a typed operation, invoke a bounded engine,
preserve provenance-bearing state, and selectively materialize the result.
Explicit tools and upfront retrieval remain general baselines and fallbacks.
Developmental pilots expose the central challenge: calculator execution and
controlled retrieval were reliable, while exposure, selection, normalization,
timing, and result use determined end-to-end behavior.  We therefore replace a
single march toward tighter coupling with capability-specific interfaces, an
eight-rung coupling ladder that includes semantic execution blocks, and
evidence gates based on resolving interface questions.  The paper contributes
a modern Transformer-era synthesis, not a claim to have invented heterogeneous
cognition, active retrieval, symbolic reasoning, or tool use.  It states a
claim hierarchy, failure taxonomy, roadmap, and broad falsification conditions
for deciding whether this organization earns its complexity.
\end{abstract}

\section{Introduction}

Transformers are powerful semantic interfaces but awkward universal
substrates.  Exact arithmetic, logical closure, constraint propagation, source
lookup, and freshness-sensitive verification already have specialized
algorithms and data structures.  Asking a language model to approximate every
such operation internally can waste capacity and conceal uncertainty.  Asking
it to acquire every fact parametrically creates a different problem: knowledge
can be absent, stale, private, source-specific, or contested.

Explicit tool calling is an important solution and should not be discarded.
Yet ``tool use'' covers interfaces with different timing, state, and control
semantics: a model can author an API call, a planner can schedule a solver, a
retriever can run before generation, a control token can request evidence, or a
runtime can recognize a bounded operation in developing text.  The systems
question is broader than whether tools should be automatic:

\begin{quote}
Which specialized computational or epistemic substrate can assist the evolving
neural computation now, through what interface, and at what coupling level?
\end{quote}

We call the design space \emph{heterogeneous cognitive Transformers}.  Here,
``cognitive'' has a narrow operational meaning concerning state participation
in ongoing inference; it makes no claim about consciousness or biological
equivalence.  The Transformer remains the general learned semantic and
generative substrate.  Specialized components recognize opportunities for
assistance, normalize typed requests, invoke bounded engines, preserve
provenance, materialize relevant outcomes, and abstain or fall back when an
automatic interface is uncertain, costly, unsafe, or unsupported.

The central thesis is architectural and falsifiable.  A useful runtime must
separate deficits, expose boundary failures, and justify tighter coupling
against strong simple alternatives.  Early program evidence is included not as
confirmation but because negative and mixed results have already changed the
proposed interfaces and evidence gates.  In particular:
\begin{itemize}
  \item computational and epistemic coprocessors form twin tracks with a
  common interrupt, typed-state, and fallback contract;
  \item different capabilities may require different interfaces rather than a
  universal automatic reflex;
  \item semantic execution blocks provide an intermediate interface between
  verbose tool calls and implicit interrupts; and
  \item later mechanisms are justified when an earlier interface question is
  sufficiently resolved, not only when every earlier mechanism wins.
\end{itemize}

\section{Two Deficits, Two Coprocessor Families}

\subsection{Competence and knowledge}

A \emph{competence deficit} occurs when the premises or operands are present,
but a specialized operation would be more reliable or efficient than neural
continuation.  Arithmetic, type checking, graph reachability, logical
inference, theorem proving, database aggregation, and bounded simulation are
examples.  A computational coprocessor returns a value, derivation,
certificate, counterexample, or explicit failure under declared semantics.

A \emph{knowledge deficit} occurs when a pending commitment needs information
that is missing, stale, source-specific, private, or insufficiently supported.
An epistemic coprocessor returns records, passages, source metadata,
timestamps, conflicts, or an explicit no-evidence status.  Retrieval provides
evidence rather than truth: a precise retrieved string may still be stale or
untrustworthy.

\begin{table}[t]
\centering
\small
\caption{Competence and knowledge deficits imply different requests, success
criteria, and characteristic risks.}
\label{tab:deficits}
\begin{tabularx}{\linewidth}{@{}p{0.16\linewidth}YYY@{}}
\toprule
 & Competence deficit & Knowledge deficit & Composed case \\
\midrule
Need & Perform or verify an operation & Acquire or verify evidence & Acquire inputs, then compute \\
Request & Typed expression, constraints, graph, program & Entity, relation, source, freshness & Evidence request plus typed operation \\
Output & Value, proof status, derivation, counterexample & Record, passage, provenance, evidence status & Provenance-linked derived result \\
Success & Formalization and execution correctness & Relevance, support, freshness, source quality & Both tests and correct composition \\
Primary risk & Wrong formalization or unsafe execution & Missing, stale, conflicting, or poisoned evidence & Laundering weak evidence through exact computation \\
\bottomrule
\end{tabularx}
\end{table}

The distinction is operational rather than ontological.  ``Which region grew
fastest?'' can require both a versioned data lookup and exact aggregation.  The
separation remains useful because it localizes error and prevents an exact
calculation from strengthening weak source provenance.

\subsection{Uncertainty and epistemic risk}

Low next-token confidence is an observable model property.  Epistemic risk is
a property of a pending commitment's dependence on external, fresh,
source-specific, or verifiable evidence.  A model can be uncertain about a
timeless paraphrase and confident about a stale fact.  Confidence can therefore
be one retrieval trigger without being a complete account of evidence need.

\section{Relationship to Adaptive Retrieval}

Retrieval-augmented generation supplies non-parametric evidence to a generator
\cite{rag}.  FLARE is a direct precedent for generation-time retrieval: it
predicts upcoming content and uses low-confidence tokens as a trigger
\cite{flare}.  Self-RAG trains the main LM to emit retrieval and reflection
control tokens concerning relevance, support, and utility \cite{selfrag}.
Adaptive-RAG, DRAGIN, corrective RAG, and multi-source retrieval further study
when to retrieve, information needs, evidence correction, and source selection
\cite{adaptiverag,dragin,crag,mspr}.  Learned adaptive retrieval, support or
relevance tokens, and generation-time retrieval are not novelty claims here.

Our narrower questions are whether semantic epistemic risk complements neural
uncertainty, whether routing should span structurally different evidence
substrates, whether control must live in the main LM, and whether a shared
runtime contract can extend from retrieval to computation.  Sidecar detectors
may be model independent; learned control in the main model may ultimately be
better.
Both are hypotheses to compare.

\begin{table}[t]
\centering
\scriptsize
\setlength{\tabcolsep}{2.5pt}
\caption{Retrieval and program interfaces.  Entries describe primary
mechanisms, not every possible extension or deployment.}
\label{tab:retrieval-program}
\begin{tabularx}{\linewidth}{@{}p{0.17\linewidth}YYYY@{}}
\toprule
Dimension & Conventional RAG & FLARE & Self-RAG & Proposed program \\
\midrule
Timing & Usually before generation & During generation & Adaptive during generation & Before or during generation \\
Control signal & Request or harness & Low confidence in predicted text & Learned retrieval/reflection tokens & Typed semantic event; later learned \\
Controller location & Orchestrator & External confidence logic & Main LM & Sidecar/runtime or main LM, compared \\
External family & Retriever & Retriever & Retriever & Compute and heterogeneous retrieval \\
State & Retrieved context & Evidence for regeneration & Evidence plus reflection signals & Typed provenance-bearing micro-state \\
Fallback & Deployment-specific & Continue/regenerate policy & Learned policy & Explicit tools and upfront RAG \\
Central question & What context helps? & When does confidence justify retrieval? & Can the LM learn retrieval/reflection? & Which substrate, interface, and coupling now? \\
\bottomrule
\end{tabularx}
\end{table}

\section{Twin Architecture and Common Contract}

Let $y_{1:t}$ denote the generated prefix and $h_t$ an optional decoding state.
A controller detects a pending need and may abstain:
\[
 d_t=D(y_{1:t},h_t,\cstate_t), \qquad
 d_t\in\{\textsc{none},\textsc{compute},\textsc{retrieve}\}.
\]
Normalization produces a typed event $e_t=N(y_{1:t},h_t,\cstate_t)$ containing
the family, operation or information need, arguments, scope, confidence, and
budget.  A router selects an engine $k=R(e_t)$ or abstains.  Execution returns
$z_t=E_k(e_t)$; the state manager records an auditable update
\[
 \cstate_{t+1}=U(\cstate_t,e_t,k,z_t,\pi_t),
\]
where $\pi_t$ captures provenance, versions, timing, and dependencies.  A
materializer exposes only relevant state as decoding continues.  Explicit tool
calls can enter this same execution and state path.

\begin{figure}[t]
\centering
\small
\stagebox{2.75in}{
\centering
\textbf{Transformer neural plane}\\[-2pt]
interpret $\rightarrow$ compose $\rightarrow$ generate\\[3pt]
$\Downarrow$ evolving text/state\\[3pt]
\textbf{Common interrupt path}\\[-2pt]
detect $\rightarrow$ expose/select $\rightarrow$ normalize typed \microir{}
$\rightarrow$ route/abstain\\[3pt]
$\swarrow\hspace{0.42\linewidth}\searrow$\\[-2pt]
\begin{tabular}{p{0.40\linewidth}p{0.40\linewidth}}
\centering\textbf{Computational family}\par
calculator, logic, constraints, graph, theorem prover, bounded code
&
\centering\textbf{Epistemic family}\par
database, lexical/vector index, knowledge graph, current web/search
\end{tabular}\\[-2pt]
$\searrow\hspace{0.42\linewidth}\swarrow$\\[-2pt]
\textbf{Common typed state path}\\[-2pt]
status, confidence, provenance, dependencies, scope, version
$\rightarrow$ selective materialization $\rightarrow$ continued decoding
}
\caption{Twin computational and epistemic planes share interrupt and state
contracts while retaining family-specific semantics and trust policies.}
\label{fig:twin-architecture}
\end{figure}

\subsection{Operational cognitive state}

We use a technical criterion rather than an anthropomorphic one:

\begin{quote}
Runtime state is treated as cognitive state when it persists beyond an isolated
service response and can participate in subsequent interpretation, reasoning,
selection, verification, or generation.
\end{quote}

A one-shot observation displayed and discarded is an ordinary tool result.
State that satisfies the criterion is typed, provenance-bearing, and
potentially persistent.  A minimal record contains a stable ID, type, payload,
engine or source version, epistemic status, confidence, scope, dependencies,
provenance, and timestamp.  Computational state should retain the formalization
and execution status, not only a number.  Epistemic state should distinguish
supported, contradicted, unverified, stale, ambiguous, and conflicting evidence.
Persistence does not imply permanence: later work must support scope,
retraction, and invalidation.

\subsection{Unified boundary failures}

Engine accuracy alone does not establish system accuracy.  Every intervention
should be attributed to a common failure taxonomy:
\begin{itemize}
  \item \textbf{Detection:} a useful interrupt is missed or a false interrupt fires.
  \item \textbf{Exposure:} the relevant operation is never represented in usable form.
  \item \textbf{Selection:} the wrong candidate, span, or subexpression is chosen.
  \item \textbf{Normalization:} intended meaning is mapped incorrectly to \microir{}.
  \item \textbf{Routing:} the wrong engine or evidence source is selected.
  \item \textbf{Execution/retrieval:} execution fails or evidence is stale,
  irrelevant, conflicting, missing, or late.
  \item \textbf{State update:} provenance, dependencies, scope, or freshness are lost.
  \item \textbf{Materialization:} too much, too little, or the wrong state is reinjected.
  \item \textbf{Use:} correct output is available but ignored or misapplied.
  \item \textbf{Override:} later generation contradicts correct active state.
  \item \textbf{Safety:} execution or an external action violates its authorization.
\end{itemize}

The difficult part of heterogeneous cognitive systems is often the
neural--coprocessor boundary, not the specialized engine itself.

\section{Capability-Specific Coupling}

\begin{quote}
Different coprocessor families may have different optimal interfaces.
\end{quote}

This is a design hypothesis, not an established result.  A calculator may
benefit from a narrow grammar or typed block; open-web search may require a
semantic and temporal trigger; consequential actions should remain explicit.
Interface choice should follow ambiguity, reversibility, latency, semantic
surface area, and harm rather than a global preference for automaticity.

\begin{table}[t]
\centering
\small
\caption{Capability-specific interface hypotheses, risks, and retained
fallbacks.}
\label{tab:capability-interfaces}
\begin{tabularx}{\linewidth}{@{}p{0.15\linewidth}p{0.23\linewidth}YY@{}}
\toprule
Capability & Likely interface & Main failure risk & Fallback \\
\midrule
Calculator & Strict grammar or semantic block & Premature or partial subexpression & Explicit calculator tool \\
Logic & Typed block/parser & Incorrect formalization & Explicit symbolic request \\
Retrieval & Epistemic interrupt & False or missed evidence need & Explicit search/upfront RAG \\
Web/current facts & Semantic plus temporal trigger & Staleness or source conflict & Explicit search with sources \\
Code & Fenced executable block & Unsafe or unbounded execution & Sandboxed tool call \\
High-risk action & Explicit invocation & Unintended side effect & User confirmation \\
\bottomrule
\end{tabularx}
\end{table}

\begin{figure}[t]
\centering
\small
\stagebox{2.35in}{
\centering
\textbf{Low ambiguity / bounded / reversible}
$\longrightarrow$
\textbf{High ambiguity / side effects / risk}\\[8pt]
\begin{tabular}{ccccc}
calculator & logic & retrieval & current web & action \\
$\downarrow$ & $\downarrow$ & $\downarrow$ & $\downarrow$ & $\downarrow$ \\
strict/block & typed block & semantic & semantic+time & explicit \\
\end{tabular}\\[10pt]
\textit{Increasing semantic breadth and risk raises the abstention and
authorization burden.}\\[5pt]
Every path may step rightward to an explicit interface; high-risk paths do not
step leftward merely because learned routing becomes available.
}
\caption{A capability-specific coupling map.  Position is a testable design
hypothesis, not a maturity ranking.}
\label{fig:coupling-map}
\end{figure}

\subsection{Semantic execution blocks}

Semantic blocks explicitly delimit a machine-executable cognitive region
without requiring the model to author a verbose API schema.  Examples are:

\begin{verbatim}
```calculator
15246377 * 746647383
```

```datalog
fact isa(penguin,bird)
fact isa(bird,animal)
query isa(penguin,animal)
```

```python
# bounded sandboxed code
```
\end{verbatim}

The runtime detects the opening block type, waits for the closing delimiter,
normalizes the enclosed content, executes it with a corresponding bounded local
engine, returns a typed result, and then continues decoding.  Waiting for the
close avoids the premature-subexpression problem of streaming lexical reflexes.
Typed blocks also expose what the model intended to offload and permit syntax,
resource, and capability checks before execution.  Arbitrary code belongs in a
no-network, resource-bounded sandbox; domain-specific engines are preferable
when they suffice.

\subsection{An eight-rung coupling ladder}

Coupling is a set of alternatives to compare, not a one-way definition of
progress.  Self-RAG-style special tokens provide a precedent for learned
control signals; semantic execution blocks generalize explicit typed control
across computational families without requiring learned tokens.

\begin{figure}[t]
\centering
\small
\stagebox{3.05in}{
\begin{tabularx}{0.91\linewidth}{@{}rY@{}}
\textbf{1} & Explicit tool or function call: universal, inspectable baseline. \\
\textbf{2} & Strict syntax/grammar reflex: narrow deterministic recognition. \\
\textbf{3} & Semantic execution block: explicit typed fenced region. \\
\textbf{4} & Lexical/syntactic semantic interrupt: runtime recognizes ordinary text. \\
\textbf{5} & Heuristic routing: transparent policy selects engines or sources. \\
\textbf{6} & Learned text-level interrupt: calibrated detector/parser with abstention. \\
\textbf{7} & Hidden-state routing: frozen neural state informs intervention. \\
\textbf{8} & Jointly trained/native interface: model and runtime co-adapt. \\
\end{tabularx}
\vspace{5pt}
\centering
more implicit and co-adapted $\longrightarrow$\\[-1pt]
\textit{Advance only for a measured boundary failure or matched benefit.}
}
\caption{Revised coupling ladder.  A higher rung is not inherently better;
cost, observability, calibration, and safety remain evaluation dimensions.}
\label{fig:coupling-ladder}
\end{figure}

The operational fallback hierarchy runs in the other direction when confidence
or safety declines: Level A uses cheap automatic reflexes; Level B uses bounded
semantic blocks; Level C uses automatic semantic or epistemic interrupts;
Level D adds learned routing; Level E uses explicit tools for arbitrary,
ambiguous, expensive, or side-effectful capabilities; and Level F requires user
confirmation for high-risk actions.  Explicit tools remain the universal slow,
general fallback.

\section{Early Evidence from the Program}

The following results are preliminary developmental evidence.  Protocols are
small, one retrieval protocol was tuned during development, and neither pilot
establishes model-family or task-family generality.  Their role here is to show
how evidence changed the architecture.

\begin{table}[t]
\centering
\scriptsize
\setlength{\tabcolsep}{2.4pt}
\caption{Developmental evidence from Papers 1 and 1.5.  Percentages are exact
accuracy unless stated otherwise; no row is a confirmatory general claim.}
\label{tab:early-evidence}
\begin{tabularx}{\linewidth}{@{}p{0.14\linewidth}p{0.32\linewidth}YY@{}}
\toprule
Study & Observed result & Boundary diagnosis & Program consequence \\
\midrule
Paper 1 strict pilot & Qwen3-0.6B XPU: LLM-only 75\%; strict reflex 75\%; explicit calculator 87.5\%; oracle 100\%. Engine exact when invoked; about 50\% of accepted candidates captured the intended full expression. & Exposure/selection and use/override failures, not calculator arithmetic, limited the reflex. & Falsifies ``automatic reflex $>$ tool call'' in this easy regime; retain explicit tools and compare other interfaces. \\
Paper 1 hard held-out extension & Same checkpoint: LLM-only 43.75\%; explicit 62.5\%; strict 37.5\%; normalized reflex 75\%; blocks 12.5\%; oracle 87.5\% on 16 arithmetic cases. Normalized vs. LLM-only exact McNemar $p=0.0625$; zero false interventions on 12 controls. & Normalization helped, but the block prompt exposed placeholder/closure failure and the oracle gap remained. Engine execution stayed exact. & Characterize and replicate the normalized interface; do not infer that blocks or automatic compute generally win. \\
Paper 1.5 retrieval pilot & Qwen3-0.6B, 12 controlled-source cases: FLARE-like confidence 83.3\% accuracy at 83.3\% retrieval; semantic OR confidence 91.7\% at 91.7\%. Upfront RAG also reached 91.7\% while retrieving all cases. & One high-confidence/high-risk stale case drove the OR difference. Confidence over-retrieved three low-confidence/low-risk cases. & Uncertainty and epistemic risk were observably distinct in one controlled case; semantic risk may complement confidence. Frozen confirmation is still required. \\
\bottomrule
\end{tabularx}
\end{table}

\subsection{What the pilots do and do not show}

In Paper 1's specified strict-reflex pilot, the calculator was exact whenever
invoked, but only roughly half of accepted candidates represented the intended
full expression, and correct outputs could be ignored or overridden.  Automatic
invocation therefore did not automatically outperform an explicit calculator.
The harder held-out extension suggests that normalization can matter, while
also showing that a superficially attractive block interface can fail if the
generation contract does not reliably close and expose executable content.

Paper 1.5 provides a more promising early signal for retrieval coupling than
the strict arithmetic reflex.  Its one distinguishing stale case supports the
possibility that semantic risk complements confidence; it does not demonstrate
general superiority over FLARE.  The protocol was tuned during development,
the source was controlled, and upfront RAG matched accuracy at greater retrieval
cost.  Negative outcomes and strong fallback performance are part of the
evidence, not inconveniences to omit.

Together, the studies motivate measuring detection, exposure, selection,
normalization, materialization, use, and override separately.  Specialized
engines can be correct while the heterogeneous system is not.

\section{Program Claims and Falsification}

\subsection{Claim hierarchy}

The program separates increasingly ambitious claims:
\begin{description}
  \item[Claim A.] Specialized engines are often more reliable or efficient for
  exact operations or external evidence.
  \item[Claim B.] Automatic generation time coupling can reduce explicit
  invocation burden in selected regimes.
  \item[Claim C.] Persistent typed cognitive state improves multi-step
  reasoning and reuse.
  \item[Claim D.] Learned semantic interrupts outperform hand-built routing.
  \item[Claim E.] Native, key/value, or latent coupling improves over textual
  or structured materialization.
  \item[Claim F.] Jointly trained heterogeneous systems develop a better
  division of labor than monolithic neural models.
\end{description}

Claim A is supported by the existence and observed correctness of bounded
engines, but does not imply Claim B.  The early pilots probe B in narrow regimes
and cannot establish C--F.  Each later claim needs its own matched baseline,
cost model, ablation, and failure analysis.

\subsection{What would falsify the broad program?}

The broad program is weakened if, across multiple domains:
\begin{itemize}
  \item explicit tool calling consistently matches or dominates automatic
  coupling on accuracy and cost;
  \item semantic interrupts create more false interventions than useful assistance;
  \item state and provenance overhead exceeds saved neural computation;
  \item heterogeneous routing provides no benefit over universal tools or retrievers;
  \item small models do not gain disproportionately;
  \item tighter interfaces do not improve scaling with task difficulty; or
  \item hidden or native coupling adds complexity without measurable value.
\end{itemize}

A successful program does not require every family to win.  A valid outcome may
be that some capabilities benefit from reflex coupling, others from execution
blocks, others from learned routing, and high-risk actions remain explicit.
Conversely, a sequence of narrow positive pilots would not validate a universal
heterogeneous architecture without cross-domain and systems-level benefit.

\section{Evidence-Gated Twin Roadmap}

The governing rule is:

\begin{quote}
A later paper may proceed when the earlier interface question has been resolved
sufficiently to motivate the next mechanism.  Do not generalize an unresolved
interface.
\end{quote}

For Paper 1 to Paper 2, a strict reflex need not win; the computational
interface must be sufficiently characterized.  A successful strict reflex can
be reused, a better semantic block can replace it, semantic-selection failure
can make learned interrupts more urgent, and persistent explicit-tool
superiority must narrow the automatic-compute claim.  For Paper 1.5 to Paper
2.5, generation-time retrieval must have value in some regime and/or oracle
source heterogeneity must show measurable headroom.  Paper 2.5 motivates Paper
3.5 only if source selection matters and heuristics leave an oracle gap.  Paper
4 and beyond require concrete provenance or retraction failures caused by
persistent multi-engine or multi-source state.

\begin{table}[p]
\centering
\scriptsize
\setlength{\tabcolsep}{2.4pt}
\caption{Twin roadmap and evidence question.  A gate resolves an interface
decision; it does not require every candidate mechanism to win.}
\label{tab:roadmap}
\begin{tabularx}{\linewidth}{@{}p{0.095\linewidth}p{0.22\linewidth}YY@{}}
\toprule
Paper & Scope & Evidence question & Decision enabled \\
\midrule
0 & Position and architecture & Are deficits, interfaces, failures, and falsifiers operational? & Freeze a testable program, not an empirical conclusion. \\
1 & Calculator/interface & Which strict, normalized, block, or explicit interface survives selection and use failures? & Choose a bounded compute interface or narrow automaticity. \\
1.5 & Single-source epistemic retrieval & Does generation-time retrieval help in any regime; does semantic risk add signal beyond confidence? & Freeze trigger questions and quantify source-heterogeneity headroom. \\
2 & Heterogeneous compute & Do multiple computational substrates add value under the characterized interface? & Keep transparent routing or identify a measured routing gap. \\
2.5 & Heterogeneous retrieval & Does source selection matter, and how large is the heuristic-to-oracle gap? & Justify or reject learned epistemic routing. \\
3 & Learned compute interrupts & Can learned detection/selection beat transparent compute policies under shift? & Justify hidden-state or joint compute coupling. \\
3.5 & Learned epistemic interrupts & Can learned trigger/source routing beat calibrated confidence and semantic heuristics? & Justify learned epistemic control. \\
4 & Transactions, provenance, retraction & Does persistent multi-family state create concrete dependency, rollback, or conflict failures? & Add only the state machinery needed by observed failures. \\
5 & Structured/PRA/native materialization & Does selective structured access improve over textual reinjection at matched cost? & Justify native or key/value interfaces. \\
6 & Co-adaptation & Does joint training improve calibrated division of labor without losing fallback robustness? & Justify end-to-end adaptation. \\
7 & Integrated runtime & Do combined capabilities improve reliability, cost, and scaling over strong modular agents? & Accept, narrow, or reject the integrated architecture. \\
\bottomrule
\end{tabularx}
\end{table}

\section{Architectural Antecedents and Synthesis}

The proposal synthesizes old and new systems ideas.  Toolformer and function
calling establish model-mediated APIs \cite{toolformer}; ReAct and
planner--executor systems interleave reasoning and action \cite{react}; PAL
delegates generated programs to interpreters \cite{pal}.  Logic-LM,
AlphaGeometry, LeanDojo, and dynamic solver composition demonstrate neural
formalization and symbolic or theorem-prover interaction
\cite{logiclm,alphageometry,leandojo,dynamic}.  Neural module networks compose
learned specialist modules \cite{nmn}.

Blackboard architectures organize heterogeneous knowledge sources around shared
state and opportunistic control \cite{blackboard}; production systems and
classical cognitive architectures provide antecedents for event-driven rules,
persistent working state, and modular cognition \cite{rete,soar,actr}.  Typed
intermediate representations such as MLIR show how explicit contracts support
heterogeneous execution \cite{mlir}; capability-restricted WebAssembly
illustrates a portable safe-execution substrate rather than a security proof by
itself \cite{wasm}.  Heterogeneous computing supplies the coprocessor analogy:
specialized substrates can coexist behind explicit scheduling, data movement,
and consistency contracts.  MemGPT and InferCept connect interruption and
persistent state to modern LM runtime management \cite{memgpt,infercept}.

\begin{table}[t]
\centering
\scriptsize
\setlength{\tabcolsep}{1.8pt}
\caption{Expanded antecedent matrix.  The proposed contribution is a unified,
evidence-gated Transformer-era synthesis rather than invention of any row.}
\label{tab:antecedents}
\begin{tabularx}{\linewidth}{@{}p{0.18\linewidth}p{0.19\linewidth}YY@{}}
\toprule
Antecedent & Primary control/state idea & What it establishes & Question retained here \\
\midrule
Tool/function calling & Model or orchestrator emits API request & General external capabilities & When should bounded needs avoid or fall back to calls? \\
ReAct/planner--executor & Interleaved model-authored plans/actions & In-trajectory external action & Can runtime recognition help without a full authored plan? \\
PAL/program-aided LMs & Generated program plus interpreter & Verifiable external execution & Which block/normalization contract is reliable? \\
FLARE/Self-RAG/adaptive retrieval & Confidence, learned tokens, or adaptive policies & Generation-time evidence control & Is semantic risk complementary; where should control live? \\
Neuro-symbolic, theorem, solver composition & Neural formalization plus formal engines & Specialized exact reasoning & How should selection and formalization failures be exposed? \\
Neural module networks & Learned composition of specialist modules & Task-dependent modular computation & Does module routing transfer to open-ended generation? \\
Blackboard/shared-state systems & Opportunistic agents over common state & Heterogeneous coordination & Which typed state and conflict semantics are necessary? \\
Production/cognitive architectures & Rules, events, modules, working memory & Persistent modular control & Which ideas survive modern Transformer evaluation? \\
Heterogeneous compute & Specialized processors plus data movement & Capability-specific substrates & What are the scheduling/materialization costs? \\
Typed IR/event-driven systems & Explicit events and contracts & Inspectable routing and lowering & How much semantics can be normalized safely? \\
Safe execution sandboxes & Capability and resource confinement & Bounded execution mechanisms & Which generated regions are safe enough to run? \\
LM runtime/memory systems & Interrupt and state scheduling & Systems costs of augmented inference & When does persistent state earn its overhead? \\
\bottomrule
\end{tabularx}
\end{table}

The synthesis differs from ``tools, but automatic.''  It treats interface type
as capability-specific, makes provenance-bearing state a first-class but gated
hypothesis, permits runtime and model control to compete, and requires automatic
paths to abstain into explicit general fallbacks.

\section{Evaluation Principles and Safety}

Every experiment should compare LLM-only, matched-prompt, explicit tool or
retrieval, automatic candidate interfaces, and an oracle boundary condition.
Report end-to-end correctness alongside detector precision/recall, candidate
coverage, selection and normalization accuracy, engine correctness, source
quality, materialization/use/override rates, latency, tokens, calls, and state
cost.  Difficulty scaling and adversarial non-trigger controls are required.
Sham interrupts can separate pause effects from result effects.  Paired tests,
confidence intervals, seeds, frozen protocols, and held-out cells should replace
claims based on aggregate point differences.

Automatic execution should be bounded, local, reversible, observable, and
side-effect free by default.  Sandboxes need explicit CPU, memory, recursion,
time, filesystem, and network policies.  Retrieved material is data, not a
control channel; prompt injection, source poisoning, stale evidence, and source
conflicts must be represented rather than silently promoted to state.  User
confirmation remains mandatory for high-risk external actions.  Provenance and
state do not guarantee truth or safety, but they make dependencies inspectable
and failures attributable.

\section{Discussion and Limitations}

The competence/knowledge distinction may blur in composed tasks, and its value
must be judged by better error localization and experimental prediction.
Interrupts can arrive too early to expose a complete operation or too late to
prevent a commitment.  Text may be a sufficient interface longer than expected;
native state access may add latency and opacity without benefit.  Persistent
state can itself become a context and consistency burden.  Small pilots on one
checkpoint cannot forecast behavior for stronger models, long-form generation,
open-web sources, or consequential domains.

Progressive Retrieval Attention is therefore not central before Paper 5.
Explicit tools, upfront RAG, transparent policies, and abstention remain valid
end states rather than temporary baselines destined for replacement.

If heterogeneous coprocessors become numerous and persistent, a later research
direction can ask whether natural language should remain the universal internal
interchange format, or whether typed/\allowbreak learned cognitive interlingua are more
appropriate.

\section{Conclusion}

A Transformer need not be either a complete monolithic reasoner or a thin tool
dispatcher.  It can remain the learned semantic and generative substrate while
bounded computational and epistemic coprocessors contribute exact operations
and external evidence.  The unifying research object is the boundary: how a
need is exposed, selected, normalized, routed, recorded, materialized, used,
and safely overridden or retracted.

Developmental Papers 1 and 1.5 sharpen rather than confirm the thesis.  They
show reliable engines behind fragile interfaces, mixed automatic results,
useful explicit fallbacks, and one controlled distinction between uncertainty
and epistemic risk.  The resulting program is capability-specific,
fallback-preserving, and evidence-gated.  It succeeds only where heterogeneous
coupling yields reproducible reliability, efficiency, scaling, or small-model
benefit over simpler alternatives; elsewhere, the correct architecture is the
simpler one.

\begin{thebibliography}{99}

\bibitem{toolformer}
T.~Schick et al.
Toolformer: Language models can teach themselves to use tools.
\emph{NeurIPS}, 2023.
\url{https://arxiv.org/abs/2302.04761}.

\bibitem{react}
S.~Yao et al.
ReAct: Synergizing reasoning and acting in language models.
\emph{ICLR}, 2023.
\url{https://arxiv.org/abs/2210.03629}.

\bibitem{pal}
L.~Gao et al.
PAL: Program-aided language models.
\emph{ICML}, 2023.
\url{https://arxiv.org/abs/2211.10435}.

\bibitem{rag}
P.~Lewis et al.
Retrieval-augmented generation for knowledge-intensive NLP tasks.
\emph{NeurIPS}, 2020.
\url{https://arxiv.org/abs/2005.11401}.

\bibitem{flare}
Z.~Jiang et al.
Active retrieval augmented generation.
\emph{EMNLP}, 2023.
\url{https://arxiv.org/abs/2305.06983}.

\bibitem{selfrag}
A.~Asai, Z.~Wu, Y.~Wang, A.~Sil, and H.~Hajishirzi.
Self-RAG: Learning to retrieve, generate, and critique through self-reflection.
\emph{ICLR}, 2024.
\url{https://arxiv.org/abs/2310.11511}.

\bibitem{adaptiverag}
S.~Jeong, J.~Baek, S.~Cho, S.~J.~Hwang, and J.~C.~Park.
Adaptive-RAG: Learning to adapt retrieval-augmented language models through
question complexity.  arXiv:2403.14403, 2024.

\bibitem{dragin}
W.~Su, Y.~Tang, Q.~Ai, Z.~Wu, and Y.~Liu.
DRAGIN: Dynamic retrieval augmented generation based on the information needs
of large language models.  arXiv:2403.10081, 2024.

\bibitem{crag}
S.-Q.~Yan, J.-C.~Gu, Y.~Zhu, and Z.-H.~Ling.
Corrective retrieval augmented generation.  arXiv:2401.15884, 2024.

\bibitem{mspr}
Q.~Zhao, R.~Wang, X.~Wang, D.~Zha, and N.~Mu.
Towards multi-source retrieval-augmented generation via synergizing reasoning
and preference-driven retrieval.  arXiv:2411.00689, 2024.

\bibitem{logiclm}
L.~Pan, A.~Albalak, X.~Wang, and W.~Y.~Wang.
Logic-LM: Empowering large language models with symbolic solvers for faithful
logical reasoning.  \emph{Findings of EMNLP}, 2023.

\bibitem{alphageometry}
T.~H.~Trinh, Y.~Wu, Q.~V.~Le, H.~He, and T.~Luong.
Solving olympiad geometry without human demonstrations.
\emph{Nature}, 2024.

\bibitem{leandojo}
K.~Yang et al.
LeanDojo: Theorem proving with retrieval-augmented language models.
\emph{NeurIPS}, 2023.

\bibitem{dynamic}
L.~Xu, P.~Beckmann, M.~Valentino, and A.~Freitas.
Adaptive LLM-symbolic reasoning via dynamic logical solver composition.
arXiv:2510.06774, 2025.

\bibitem{nmn}
J.~Andreas, M.~Rohrbach, T.~Darrell, and D.~Klein.
Neural module networks.  \emph{CVPR}, 2016.

\bibitem{blackboard}
B.~Hayes-Roth.
A blackboard architecture for control.
\emph{Artificial Intelligence}, 26(3):251--321, 1985.

\bibitem{rete}
C.~L.~Forgy.
Rete: A fast algorithm for the many pattern/many object pattern match problem.
\emph{Artificial Intelligence}, 19(1):17--37, 1982.

\bibitem{soar}
A.~Newell.
\emph{Unified Theories of Cognition}.
Harvard University Press, 1990.

\bibitem{actr}
J.~R.~Anderson et al.
An integrated theory of the mind.
\emph{Psychological Review}, 111(4):1036--1060, 2004.

\bibitem{mlir}
C.~Lattner et al.
MLIR: Scaling compiler infrastructure for domain specific computation.
\emph{CGO}, 2021.

\bibitem{wasm}
A.~Haas et al.
Bringing the Web up to speed with WebAssembly.
\emph{PLDI}, 2017.

\bibitem{memgpt}
C.~Packer et al.
MemGPT: Towards LLMs as operating systems.
arXiv:2310.08560, 2023.

\bibitem{infercept}
R.~Abhyankar, Z.~He, V.~Srivatsa, H.~Zhang, and Y.~Zhang.
InferCept: Efficient intercept support for augmented large language model
inference.  arXiv:2402.01869, 2024.

\end{thebibliography}

\end{document}
